\nuevaReceta{Pelotas Almerienses}{6 personas}{3 h}{receta:pelotas_almerienses}

    \begin{minipage}[t]{0.35\textwidth}
        \small
        \section*{Ingredientes}
        \begin{itemize}[leftmargin=*, itemsep=1pt]
            \item 1 Perdiz limpia
            \item 1/2 Gallina
            \item 1/2 Pollo
            \item 250 g de carne picada de cerdo
            \item 100 g de jamón serrano
            \item 150 g de miga de pan
            \item 3 Huevos
            \item 2 Dientes de ajo
            \item Perejil fresco
            \item 1 Hueso de jamón
            \item 2 Zanahorias
            \item 1 Puerro
            \item 1 Rama de apio
            \item Azafrán o colorante
            \item Pimienta negra
            \item Sal
        \end{itemize}
    \end{minipage}
    \hfill
    \begin{minipage}[t]{0.6\textwidth}
        \small
        \section*{Preparación}
        \begin{enumerate}[itemsep=2pt, topsep=2pt]
            \item \textbf{Preparar el caldo:} Colocar la perdiz, la gallina, el pollo y el hueso de jamón en una olla grande. Cubrir con agua fría y llevar lentamente a ebullición.
            \item \textbf{Limpiar y cocer:} Retirar la espuma de la superficie. Añadir las zanahorias, el puerro, el apio y una pizca de sal. Cocer a fuego suave durante unas 2 horas, hasta que las carnes estén tiernas.
            \item \textbf{Preparar las carnes:} Sacar la perdiz, la gallina y el pollo. Dejar que se templen, retirar la piel y los huesos, y picar finamente una parte de cada carne. Reservar el resto para servir con el caldo.
            \item \textbf{Hacer la masa:} Mezclar las carnes picadas con la carne de cerdo, el jamón, la miga de pan, los huevos, los ajos y el perejil muy picados. Añadir pimienta y ajustar de sal.
            \item \textbf{Formar las pelotas:} Trabajar la mezcla hasta que quede ligada. Formar pelotas de tamaño mediano con las manos ligeramente humedecidas.
            \item \textbf{Colar el caldo:} Colar el caldo, devolverlo a la olla y añadir unas hebras de azafrán o un poco de colorante.
            \item \textbf{Cocer las pelotas:} Cuando el caldo hierva suavemente, incorporar las pelotas con cuidado. Cocer durante 20-25 minutos, sin remover bruscamente para que no se rompan.
            \item \textbf{Servir:} Repartir el caldo bien caliente con las pelotas y unos trozos de perdiz, gallina y pollo reservados.
        \end{enumerate}
    \end{minipage}

    \vspace{0.5em}
    \begin{tcolorbox}[colback=orange!5, colframe=orange!50, title=Consejo, size=small]
        Prepara el caldo el día anterior: reposado gana sabor y resulta más sencillo retirar la grasa. Si la masa queda blanda, añade un poco más de miga de pan antes de formar las pelotas.
    \end{tcolorbox}
\end{receta}
